\section*{Chapitre \ref{ch:mx:the-limit-order-book} --- Le carnet d'ordres à cours limité}

\begin{solution}{exo:mx:the-limit-order-book:1}
La vente s'exécute contre le niveau dans l'ordre temporel~: 300 actions contre le premier ordre et 150 contre le second. Le meilleur cours acheteur reste à 100.00 avec
$50+200=250$ actions réparties sur deux ordres~; le meilleur cours vendeur est inchangé à 100.01.
\end{solution}

\begin{solution}{exo:mx:the-limit-order-book:2}
$P_3=\tfrac{1}{1.2}\cdot\tfrac{1.1}{1.3}\cdot\tfrac{1.2}{1.4}\cdot\tfrac{1.3}{1.5}=0.524$.
\end{solution}

\begin{solution}{exo:mx:the-limit-order-book:3}
La réduction conserve la place de l'ordre (un message X de 200 actions dans un flux de type ITCH)~; l'augmentation modifie sa priorité sur CME Globex, qui la reporte derrière tous les ordres à son prix.
Sur Nasdaq, l'augmentation est un remplacement annulation-ajout~: un message Order Replace supprime l'ancienne référence et ajoute une nouvelle
référence à la fin de la file.
\end{solution}

\begin{solution}{exo:mx:the-limit-order-book:4}
De 99.00 à 101.00 en centimes~: $2\times100+1=201$ cases. Un mouvement de 3~\% sort le meilleur cours de la bande~: le carnet doit se recentrer (reconstruire le tableau
autour du nouveau prix) ou conserver les ordres situés hors de la bande dans une table de débordement, et le recentrage crée un pic de latence précisément au moment où le marché bouge.
\end{solution}

\begin{solution}{exo:mx:the-limit-order-book:5}
$1/0.5+\sum_{k=1}^{10}1/(0.5+0.02k)=2+16.54=18.5$ secondes.
\end{solution}

\begin{solution}{exo:mx:the-limit-order-book:6}
Le meilleur niveau est celui où toutes les exécutions ont lieu, et c'est aussi là que les ordres périmés sont retirés en premier lorsque le prix le traverse~; à un tick
derrière, les ordres s'accumulent hors de portée des petits \omterm{def:mx:the-limit-order-book:orders}{ordres au marché} et sont alimentés par ceux qui y arrivent, si bien qu'ils s'entassent. La forme est la bosse
que le chapitre~3 retrouve dans les carnets réels.
\end{solution}

\begin{solution}{exo:mx:the-limit-order-book:7}
Avec des \omterm{def:mx:the-limit-order-book:orders}{ordres au marché} au taux 1, la fourchette passe de 2.55 à 1.64 tick et le meilleur niveau compte 3.29 ordres au lieu de 2.72~: les \omterm{def:mx:the-limit-order-book:orders}{ordres au marché} vident
le meilleur niveau deux fois moins souvent. Derrière cinq ordres, la probabilité d'exécution vaut 0.602 contre 0.769 pour la proposition, soit une perte de 22~\% (23~\% au taux 2)~: les attentes
sont deux fois plus longues, ce qui laisse aux ordres suivants davantage de temps pour améliorer le prix, mais une fourchette plus étroite laisse moins de prix à améliorer, et les deux
effets se compensent presque.
\end{solution}

\begin{solution}{exo:mx:the-limit-order-book:8}
Le facteur $\mu/(\mu+\nu)$ suppose que personne ne peut prendre la place de l'ordre en tête. Tout cours acheteur placé au-dessus devient le meilleur et reçoit d'abord les ventes au marché~;
le niveau peut cesser d'être le meilleur, et l'ordre attend alors indéfiniment. Dix secondes en tête renseignent aussi sur les taux~:
sur un marché calme, $\mu$ est inférieur à sa moyenne.
\end{solution}

\begin{solution}{pb:mx:the-limit-order-book:1}
\textbf{1.} 31\,535 messages~: 49.7~\% d'ajouts, 47.8~\% d'annulations, 2.5~\% d'exécutions.
\textbf{2.} 40 messages par transaction~; 3.8~\% des ordres sortis du carnet l'ont été par une transaction.
\textbf{3.} 15.5 secondes en médiane~; 24.9 pour les ordres exécutés, 15.3 pour ceux annulés~: être exécuté prend plus de temps que renoncer.
\textbf{4.} 4.5 secondes pour le meilleur cours acheteur (4.8 pour le meilleur cours vendeur).
\textbf{5.} 2\,834 ordres, avec une médiane de 1\,400 actions devant.
\textbf{6.} 13~\% sont exécutés au moins en partie, 10~\% en totalité.
\textbf{7.} Un tick derrière le meilleur (1\,508 actions du côté acheteur contre 1\,366 au meilleur)~: le meilleur niveau est vidé par les exécutions et les
annulations d'ordres périmés, le suivant ne l'est pas.
\textbf{8.} La durée de vie des ordres et la part des ordres exécutés augmenteraient~; les files devant seraient plus longues et la profondeur plus grande~; le nombre de messages par transaction baisserait.
\textbf{9.} Ordres unitaires~; \omterm{def:mx:the-limit-order-book:orders}{ordres au marché} au taux $\mu$ depuis la tête~; annulations au taux $\theta$ par ordre devant et $\nu$ pour le nôtre~; horloges exponentielles
indépendantes~; personne ne passe devant. Alors $P_n=\tfrac{\mu}{\mu+\nu}\prod_{k=1}^{n}\tfrac{\mu+k\theta}{\mu+k\theta+\nu}$.
\textbf{10.} 0.952, 0.870 et 0.755.
\textbf{11.} \texttt{queue\_mc(5, 2, 0.05, 0.05)} avec 20\,000 trajectoires concorde à 0.01 près (0.869).
\textbf{12.} 2.83 secondes.
\textbf{13.} 0.953 derrière un ordre, 0.871 derrière cinq~: la file du meilleur cours acheteur se comporte alors exactement comme le suppose le modèle, car personne ne peut placer
un cours acheteur au-dessus d'elle.
\textbf{14.} 0.912, 0.672 et 0.535.
\textbf{15.} \emph{Résultat notable}~: la probabilité d'exécution derrière $n$ ordres passe de 0.95 ($n=1$) à 0.76 ($n=12$) dans la proposition~; dans le marché
à intelligence nulle avec amélioration de prix, elle passe de 0.91 à 0.54 derrière huit ordres, et l'amélioration retire 23~\% des exécutions promises derrière
cinq et 35~\% derrière huit.
\textbf{16.} Seuls les \omterm{def:mx:the-limit-order-book:orders}{ordres au marché} et les annulations retirent des ordres des meilleurs niveaux, et rien n'est jamais placé à l'intérieur de la fourchette, si bien que le meilleur prix de chaque côté
ne fait que s'éloigner.
\textbf{17.} 2.55 ticks et 2.72 ordres.
\textbf{18.} Personne ne passe devant (amélioration de prix), et les taux constants~: les deux biais poussent la proposition à la hausse~; dans un carnet réel, les flux au meilleur niveau
se regroupent aussi dans le temps (le livre~4, chapitre~7).
\textbf{19.} $\mu$ à partir des exécutions au meilleur niveau par unité de temps, $\theta$ à partir des annulations au meilleur niveau divisées par la taille de la file intégrée dans le temps,
$\nu$ à partir de la durée de vie de ses propres ordres~; le tout conditionnellement à l'état de la file.
\textbf{20.} La probabilité, et le temps, d'être exécuté avant que le prix ne bouge~: la tête d'une longue file vaut une exécution, sa queue une loterie.
\end{solution}

\begin{solution}{iq:mx:the-limit-order-book:1}
Une table de hachage de la référence d'ordre vers l'ordre (annulations et exécutions la désignent), une table triée ou un tableau indexé par le prix vers le niveau, et une liste FIFO
par niveau. Annulation et exécution~: $O(1)$ avec une liste intrusive~; ajout~: $O(1)$, ou $O(\log L)$ lorsqu'un arbre reçoit un nouveau niveau~; meilleur prix~: $O(1)$ avec un
pointeur mis en cache. \iqlookfor{les trois chemins d'accès, la table de hachage pour les références, pourquoi les tableaux l'emportent lorsque les prix se regroupent.}
\end{solution}

\begin{solution}{iq:mx:the-limit-order-book:2}
Une diminution ne retire que les actions de l'ordre lui-même, ce qui ne nuit à personne derrière~; une augmentation permettrait à un trader de tenir la tête avec un petit ordre et
de l'agrandir lorsqu'une grosse contrepartie apparaît, en passant devant tous ceux qui attendaient. \iqlookfor{l'équité envers les ordres qui suivent, et les manipulations qu'une hausse de taille
permettrait.}
\end{solution}

\begin{solution}{iq:mx:the-limit-order-book:3}
Avec quatre ordres devant~: $P=\tfrac{\mu}{\mu+\nu}\prod_{k=1}^{4}\tfrac{\mu+k\theta}{\mu+k\theta+\nu}$, issue d'horloges exponentielles concurrentes et de l'absence de mémoire.
\iqlookfor{la décomposition en étapes et l'indépendance des étapes~; une remarque selon laquelle l'amélioration de prix la fait baisser.}
\end{solution}

\begin{solution}{iq:mx:the-limit-order-book:4}
Lorsque le prix sort de la bande que couvre le tableau (un marché rapide, une reprise après suspension très éloignée), ou pour des instruments aux amplitudes de cours énormes et aux carnets clairsemés.
Recentrer le tableau, conserver une table de débordement pour les ordres éloignés et dimensionner la bande d'après la volatilité de l'instrument. \iqlookfor{le compromis entre mémoire et amplitude,
et le pic de latence du recentrage au pire moment.}
\end{solution}

\begin{solution}{iq:mx:the-limit-order-book:5}
Pas en soi~: les annulations sont la manière dont les fournisseurs de liquidité mettent à jour leurs cotations lorsqu'une information arrive, et le carnet affiché reste à jour. Elles coûtent
de la capacité de messages et rendent la profondeur affichée moins ferme (l'évanouissement de la cotation, chapitre~18)~; les régulateurs les surveillent au moyen des ratios ordres sur transactions (chapitre~24).
\iqlookfor{une réponse dans les deux sens~: les cotations périmées sont pires, mais la liquidité éphémère induit en erreur.}
\end{solution}

\begin{solution}{iq:mx:the-limit-order-book:6}
Juste~: l'ordre de grandeur de la fourchette et sa dépendance aux taux de flux d'ordres, la forme de la profondeur près du meilleur niveau, la raison pour laquelle la plupart des ordres ne sont pas exécutés.
Faux~: le flux d'ordres n'est pas indépendant (le signe des transactions a une mémoire longue, chapitre~12), les taux réagissent à l'état du carnet et à l'information, et les prix n'ont
aucun ancrage dans une valeur. \iqlookfor{c'est un \omterm{def:mx:the-limit-order-book:zi}{modèle à intelligence nulle}~: une référence pour ce qui s'explique sans intelligence.}
\end{solution}
